\section*{अध्याय \ref{ch:g6:fractions} --- भिन्नें: पहले कदम}

\begin{solution}{exo:g6:fractions:1}
केक: $\frac38$ खाया गया। \quad चॉकलेट: $\frac{7}{12}$। \quad
पाई: $\frac66 = 1$ (पूरी पाई)।
\end{solution}

\begin{solution}{exo:g6:fractions:2}
पाँच में से चार हिस्से रँगे जाते हैं। $\frac45 < 1$ (पाँच में से
चार हिस्से पूरी पट्टी से कम हैं); $\frac65 > 1$ (छह पाँचवें भाग
एक पूरी पट्टी और एक पाँचवाँ भाग हैं)।
\end{solution}

\begin{solution}{exo:g6:fractions:3}
तिहाइयों में अंकित रेखा पर: $\frac13$, $0$ के एक निशान बाद है;
$\frac53$, $2$ से एक निशान पहले (क्योंकि $2 = \frac63$);
$\frac73$, $2$ के एक निशान बाद। तो $\frac53$ और $\frac73$ दोनों
$2$ से $\frac13$ की दूरी पर हैं।
\end{solution}

\begin{solution}{exo:g6:fractions:4}
$\frac94 = 2 + \frac14$; \quad
$\frac{13}{5} = 2 + \frac35$; \quad
$\frac{12}{6} = 2$ (ठीक-ठीक); \quad
$\frac{25}{8} = 3 + \frac18$।
\end{solution}

\begin{solution}{exo:g6:fractions:5}
$\frac23 = \frac{8}{12}$ ($4$ से गुणा); \quad
$\frac{15}{20} = \frac34$ ($5$ का भाग); \quad
$\frac{7}{10} = \frac{21}{30}$ ($3$ से गुणा); \quad
$\frac56 = \frac{10}{12}$ ($2$ का भाग)।
\end{solution}

\begin{solution}{exo:g6:fractions:6}
$\frac68 = \frac34$; \quad
$\frac{15}{25} = \frac35$; \quad
$\frac{18}{24} = \frac{9}{12} = \frac34$; \quad
$\frac{35}{7} = 5$ (एक पूर्ण संख्या)।
\end{solution}

\begin{solution}{exo:g6:fractions:7}
$86$ का $\frac12$: $86 \div 2 = 43$।

$60$ का $\frac34$: $60 \div 4 = 15$, फिर $15 \times 3 = 45$।

$35$ का $\frac25$: $35 \div 5 = 7$, फिर $7 \times 2 = 14$।

$250$ का $\frac{7}{10}$: $250 \div 10 = 25$, फिर $25 \times 7 = 175$।
\end{solution}

\begin{solution}{exo:g6:fractions:8}
$\frac12 > \frac13$: आधे, तिहाइयों से बड़े हिस्से होते हैं।
$\frac35 > \frac25$: वही हिस्से, पर ज़्यादा।
$\frac34 = \frac68 < \frac78$: सात आठवें भाग, छह आठवें भागों पर भारी।
\end{solution}

\begin{solution}{exo:g6:fractions:9}
$\frac12 = 0.5$; \quad $\frac34 = 0.75$; \quad
$\frac{7}{10} = 0.7$; \quad $\frac94 = 2.25$। अलग पड़ती है
$\frac13 = 0.333\dots$, जिसका दशमलव लेखन कभी ख़त्म नहीं होता: वह
\omterm{def:g6:decimals:places}{दशमलव संख्या} नहीं है।
\end{solution}

\begin{solution}{exo:g6:fractions:10}
बस से आने वाले: $28$ का $\frac34$: $28 \div 4 = 7$, फिर
$7 \times 3 = 21$ विद्यार्थी। पैदल आने वाले: $28 - 21 = 7$
विद्यार्थी (यही बची हुई \omterm{def:g3:division:half}{चौथाई} है: $28 \div 4 = 7$)। जाँच: $21 + 7 = 28$।
\end{solution}

\begin{solution}{exo:g6:fractions:11}
बचत के दो तिहाई $18$ यूरो हैं, इसलिए \emph{एक} तिहाई
$18 \div 2 = 9$ यूरो है, और पूरी बचत (तीन तिहाई)
$9 \times 3 = 27$ यूरो थी। जाँच: $27$ का $\frac23$, $18$ है।
\end{solution}

\begin{solution}{exo:g6:fractions:12}
इस्तेमाल किया गया पानी टंकी के $\frac58$ से $\frac48$ तक ले आया
(\omterm{def:g3:division:half}{आधा} यानी चार आठवें भाग): $10$ लीटर \omterm{def:g2:measure:units}{धारिता} के
$\frac58 - \frac48 = \frac18$ हैं। तो पूरी टंकी में
$10 \times 8 = 80$ लीटर आते हैं। जाँच: $80$ का $\frac58$, $50$ है;
$10$ घटाने पर $40$ बचता है, जो सचमुच $80$ का \omterm{def:g3:division:half}{आधा} है।
\end{solution}

\begin{solution}{pb:g6:fractions:1}
\textbf{1.} $16$ ख़ानों का \omterm{def:g3:division:half}{आधा} $8$ ख़ाने है। \omterm{def:g6:fractions:def}{भिन्न} के रूप में:
चॉकलेट का $\frac12$, और चूँकि हर \omterm{def:g3:division:half}{आधा} $16$ में से $8$ ख़ाने है,
इसलिए $\frac{8}{16}$ भी --- वही \omterm{def:g6:fractions:def}{भिन्न} दो तरह लिखी हुई
(\cref{prop:g6:fractions:equal})।

\textbf{2.} बचे हुए $8$ ख़ानों का \omterm{def:g3:division:half}{आधा} $4$ ख़ाने है, जो चॉकलेट का
$\frac{4}{16}$ है। $4$ से सरल करने पर:
$\frac{4}{16} = \frac14$। चित्र पर: बचे हुए आधे को दो में काटने
से मूल चॉकलेट की दो चौथाइयाँ बनती हैं --- आधे का \omterm{def:g3:division:half}{आधा} एक
\omterm{def:g3:division:half}{चौथाई} है।

\textbf{3.} तीसरा कौर: $4$ ख़ानों का \omterm{def:g3:division:half}{आधा} $2$ ख़ाने है, यानी
चॉकलेट का $\frac{2}{16} = \frac18$। चौथा कौर: $1$ ख़ाना,
यानी $\frac{1}{16}$। एक कौर से अगले तक हर
\emph{दुगुना} होता जाता है: $2$, $4$, $8$, $16$।

\textbf{4.} खाए गए: $16$ में से $8 + 4 + 2 + 1 = 15$ ख़ाने, यानी
चॉकलेट का $\frac{15}{16}$। बचा: $1$ ख़ाना, $\frac{1}{16}$।

\textbf{5.} $\frac12 = \frac{8}{16}$ (ऊपर और नीचे $8$ से गुणा),
$\frac14 = \frac{4}{16}$ ($4$ से), $\frac18 = \frac{2}{16}$
($2$ से), और $\frac{1}{16}$ वैसी ही रहती है। सोलहवें भाग गिनो:
वे $8 + 4 + 2 + 1 = 15$ हैं, इसलिए कुल $\frac{15}{16}$ है।

\textbf{6.} कौर: $32$, $16$, $8$, $4$, $2$, $1$ ख़ाने। कुल खाया
गया: $32 + 16 + 8 + 4 + 2 + 1 = 63$ ख़ाने, इसलिए $1$ ख़ाना
बचता है: चॉकलेट का $\frac{1}{64}$।

\textbf{7.} हर कौर के बाद बचता है
\[
\frac12,\quad \frac14,\quad \frac18,\quad \frac{1}{16},\quad
\frac{1}{32},\quad \frac{1}{64} :
\]
\omterm{def:g3:division:remainder}{शेषफल} हर बार \omterm{def:g3:division:half}{आधा} हो जाता है --- उसका हर दुगुना होता जाता है।
सातवाँ, काल्पनिक कौर $\frac{1}{128}$ छोड़ेगा।

\textbf{8.} हर कौर बचे हुए का केवल \emph{\omterm{def:g3:division:half}{आधा}} खाता है, इसलिए
\omterm{def:g3:division:remainder}{शेषफल} का दूसरा \omterm{def:g3:division:half}{आधा} वहीं रह जाता है: कितने भी कौरों के बाद कुछ न
कुछ हमेशा बचा रहता है। इसलिए खाया हुआ कुल हमेशा $1$ से नीचे रहता
है --- ठीक \cref{pb:g6:decimals:1} के बीस नौ की तरह, जो $3$ के
और-और पास आते हैं पर उस तक पहुँचते नहीं।

\textbf{9.} दोनों भिन्नों का \omterm{def:g4:fractions:def}{अंश} $1$ है, और $128$ हिस्सों में
काटने से $100$ हिस्सों में काटने की तुलना में छोटे हिस्से बनते
हैं: इसलिए $\frac{1}{128} < \frac{1}{100}$। सातवें कौर के बाद
\omterm{def:g3:division:remainder}{शेषफल} $\frac{1}{128}$ है, इसलिए खाया हुआ भाग
$1 - \frac{1}{100} = \frac{99}{100}$ से आगे निकल जाता है:
चुनौती सातवें कौर पर जीत ली जाती है।

\textbf{10.} सोलहवें भागों में कुल हैं $\frac{8}{16}$,
$\frac{12}{16}$, $\frac{14}{16}$, $\frac{15}{16}$। हर नया कुल
ठीक पिछले कुल और $1$ के \emph{बीचों-बीच} गिरता है: सीढ़ी $1$ तक
बची अपनी दूरी को हर बार \omterm{def:g3:division:half}{आधा} करती जाती है, पर उस पर कभी पैर नहीं रखती।

\textbf{11.} सही क्या है: ज़ेनो के वर्णन के हर चरण पर सचमुच कुछ
दूरी बची रहती है (प्रश्न 8) --- चरणों की सूची कभी ख़त्म नहीं
होती। जाल यह है: चरण बेहद छोटे होते जाते हैं, दूरी में
\emph{और} दौड़ने के समय में भी; धावक हर चरण पर बराबर समय नहीं
लगाता। दौड़ को अनगिनत सिकुड़ते टुकड़ों में बताने से दौड़ ख़ुद
अनंत नहीं हो जाती: धावक दीवार तक पहुँच जाता है, और प्रश्न 9
दिखाता है कि वर्णन दीवार से पहले के हर बिंदु को पार कर जाता है।

\textbf{12.} $8$ चॉकलेटें आधों में काटो: $16$ आधे टुकड़े। $4$
चॉकलेटें चौथाइयों में: $16$ \omterm{def:g3:division:half}{चौथाई} टुकड़े। $2$ चॉकलेटें आठवें
भागों में: $16$ टुकड़े। आख़िरी चॉकलेट सोलहवें भागों में:
$16$ छोटे टुकड़े। इसमें $8 + 4 + 2 + 1 = 15$ चॉकलेटें लगती हैं,
और हर बच्चे को हर नाप का एक-एक टुकड़ा मिलता है।

\textbf{13.} हर बच्चे का हिस्सा है
\[
\frac12 + \frac14 + \frac18 + \frac{1}{16} = \frac{15}{16}
\]
चॉकलेट का (प्रश्न 5)। ऐसे सोलह हिस्से मिलकर
$16 \times \frac{15}{16} = 15$ चॉकलेटें बनाते हैं
(\cref{thm:g6:fractions:quotient}): ठीक उतनी ही जितनी काटी गई
थीं, कुछ बचे बिना --- केवल आधा-आधा करके $15$ चॉकलेटों का $16$
बच्चों में सही बँटवारा।

\textbf{14.} पड़ोसियों को हर एक को चॉकलेट का $\frac78$ मिलता है।
हर $16$ के साथ: $\frac78 = \frac{14}{16}$, जबकि हमारे बच्चों को
$\frac{15}{16}$ मिलता है। तो $16$ वाले बच्चे को ज़्यादा मिलता है:
$\frac{15}{16} > \frac{14}{16}$, चॉकलेट के एक सोलहवें भाग से।

\textbf{15.} वर्ग वाले चित्र में हर रँगा हुआ कौर अब तक बिना रँगे
कोने का \omterm{def:g3:division:half}{आधा} भर देता है, और बिना रँगा कोना सिकुड़ता जाता है: हर
कौर के बाद वह पहले से \omterm{def:g3:division:half}{आधा} रह जाता है, और कभी ग़ायब नहीं होता।
दोनों सच्चे तथ्य: (i) सीमित संख्या में कौरों के बाद हमेशा एक
\omterm{def:g3:division:remainder}{शेषफल} बचता है --- खाया हुआ कुल हमेशा $1$ से ठीक नीचे रहता है;
(ii) पर्याप्त कौर लेने पर वह \omterm{def:g3:division:remainder}{शेषफल} किसी भी बताई गई \omterm{def:g6:fractions:def}{भिन्न}
($\frac{1}{100}$, $\frac{1}{1000}$, \dots) से छोटा हो जाता है।
``अनंत \omterm{def:g1:addition:def}{योगफल} $1$ के बराबर है'' --- इन शब्दों को ठीक-ठीक अर्थ देना
\emph{सीमाओं} का काम है, जो अगले खंड में आता है।

\end{solution}
